% ============================================================
% Section 190: 一维晶体的能带、有效质量与电场输运
% ============================================================
\section{固体理论：一维晶体的能带、有效质量与电场输运}
\label{sec:1d-band-transport}

\begin{modelbox}[题目：一维紧束缚-like能带的完整分析]
设一维晶体的电子能带可写成
\begin{equation}
E(k)=\frac{\hbar^2}{ma^2}\Bigl(\frac{7}{8}-\cos ka+\frac{1}{8}\cos 2ka\Bigr),
\end{equation}
其中 $a$ 是晶格常数。试求：
\begin{enumerate}
    \item 能带的宽度；
    \item 电子在波矢 $k$ 的状态时的速度；
    \item 能带底部和顶部电子的有效质量；
    \item 若晶格常数 $a=2.5\,\text{\AA}$，当外加 $10^3\,\text{V/m}$ 和 $10^7\,\text{V/m}$ 电场时，试分别估计电子自能带底运动到能带顶所需的时间。
\end{enumerate}
\end{modelbox}

\begin{tipbox}[考点快速分析]
\begin{itemize}
    \item \textbf{能带宽度}：求 $E(k)$ 在 $[-\pi/a,\pi/a]$ 上的最大值与最小值之差
    \item \textbf{电子群速度}：$v(k)=\frac{1}{\hbar}\frac{\dd E}{\dd k}$
    \item \textbf{有效质量}：$m^*=\hbar^2\bigl(\frac{\dd^2E}{\dd k^2}\bigr)^{-1}$
    \item \textbf{Bloch 振荡时间}：$\hbar\frac{\dd k}{\dd t}=-eE$（或 $eE$，取决于电子/空穴），$\Delta t=\frac{\hbar\Delta k}{eE}$
\end{itemize}
\end{tipbox}

\begin{derivationbox}[标准解题过程]
\textbf{(1) 能带宽度}

求极值点：
\begin{equation}
\frac{\dd E}{\dd k}=\frac{\hbar^2}{ma^2}\Bigl(a\sin ka-\frac{a}{4}\sin 2ka\Bigr)
=\frac{\hbar^2}{ma}\sin ka\Bigl(1-\frac{1}{2}\cos ka\Bigr)=0.
\end{equation}

解得 $\sin ka=0$，即 $ka=0,\pm\pi$（在第一布里渊区内）。

\begin{itemize}
    \item 能带底部 $k=0$：
    \begin{equation}
    E(0)=\frac{\hbar^2}{ma^2}\Bigl(\frac{7}{8}-1+\frac{1}{8}\Bigr)=0.
    \end{equation}
    \item 能带顶部 $k=\pm\pi/a$：
    \begin{equation}
    E\Bigl(\pm\frac{\pi}{a}\Bigr)=\frac{\hbar^2}{ma^2}\Bigl(\frac{7}{8}+1+\frac{1}{8}\Bigr)=\frac{2\hbar^2}{ma^2}.
    \end{equation}
\end{itemize}

能带宽度
\begin{equation}
\boxed{\Delta E=E_{\max}-E_{\min}=\frac{2\hbar^2}{ma^2}.}
\end{equation}

\textbf{(2) 电子速度}

\begin{equation}
\boxed{v(k)=\frac{1}{\hbar}\frac{\dd E}{\dd k}=\frac{\hbar}{ma}\Bigl(\sin ka-\frac{1}{4}\sin 2ka\Bigr).}
\end{equation}

\textbf{(3) 有效质量}

\begin{equation}
\frac{\dd^2E}{\dd k^2}=\frac{\hbar^2}{m}\Bigl(\cos ka-\frac{1}{2}\cos 2ka\Bigr).
\end{equation}

能带底部 $k=0$：
\begin{equation}
\frac{\dd^2E}{\dd k^2}\bigg|_{k=0}=\frac{\hbar^2}{m}\Bigl(1-\frac{1}{2}\Bigr)=\frac{\hbar^2}{2m},
\quad\Rightarrow\quad
\boxed{m^*_{\text{底}}=2m.}
\end{equation}

能带顶部 $k=\pm\pi/a$：$\cos\pi=-1$，$\cos 2\pi=1$，
\begin{equation}
\frac{\dd^2E}{\dd k^2}\bigg|_{k=\pi/a}=\frac{\hbar^2}{m}\Bigl(-1-\frac{1}{2}\Bigr)=-\frac{3\hbar^2}{2m},
\quad\Rightarrow\quad
\boxed{m^*_{\text{顶}}=-\frac{2m}{3}.}
\end{equation}

\textbf{(4) 电场下运动时间}

半经典运动方程：$\hbar\frac{\dd k}{\dd t}=-eE$（电子带负电，电场力与 $E$ 反向）。

从能带底 $k=0$ 到能带顶 $k=\pi/a$，$\Delta k=\pi/a$：
\begin{equation}
\Delta t=\frac{\hbar\Delta k}{eE}=\frac{\hbar\pi}{aeE}.
\end{equation}

代入 $a=2.5\,\text{\AA}=2.5\times10^{-10}\,\text{m}$，$\hbar=1.055\times10^{-34}\,\text{J}\cdot\text{s}$：
\begin{equation}
\hbar\pi/a\approx\frac{1.055\times10^{-34}\times\pi}{2.5\times10^{-10}}
\approx1.33\times10^{-24}\,\text{kg}\cdot\text{m/s}.
\end{equation}

\begin{itemize}
    \item $E=10^3\,\text{V/m}$：
    \begin{equation}
    \Delta t=\frac{1.33\times10^{-24}}{1.6\times10^{-19}\times10^3}
    \approx\boxed{8.3\times10^{-9}\,\text{s}=8.3\,\text{ns}.}
    \end{equation}
    \item $E=10^7\,\text{V/m}$：
    \begin{equation}
    \Delta t=\frac{1.33\times10^{-24}}{1.6\times10^{-19}\times10^7}
    \approx\boxed{8.3\times10^{-13}\,\text{s}=0.83\,\text{ps}.}
    \end{equation}
\end{itemize}
\end{derivationbox}

\begin{insightbox}[物理图像]
\begin{itemize}
    \item 有效质量 $m^*$ 的符号决定载流子类型：$m^*>0$ 对应能带底（电子），$m^*<0$ 对应能带顶（空穴）。
    \item 电子在恒定电场下的 $k$ 空间运动是匀速的：$k(t)=k_0-eEt/\hbar$。当电子到达布里渊区边界时，通过 Umklapp 过程（与晶格动量交换 $2\pi/a$）回到另一边界，形成 \textbf{Bloch 振荡}。
    \item 实际晶体中，电子在到达边界前就会因散射（声子、杂质）而弛豫，因此通常观察不到完整的 Bloch 振荡。只有在超晶格等人工周期结构中，由于晶格常数大、能带窄，才有可能实现。
\end{itemize}
\end{insightbox}

\begin{warnbox}[常见失误]
\begin{itemize}
    \item 有效质量公式分母是 $\frac{\dd^2E}{\dd k^2}$，不是 $\frac{\dd^2E}{\dd k^2}$ 的绝对值。顶部曲率为负，有效质量自然为负。
    \item Bloch 振荡时间公式中电子电荷取 $e>0$，但电子受力为 $-eE$，故 $\Delta k$ 从 $0$ 到 $\pi/a$ 时 $\Delta t$ 为正。
    \item 注意区分 $E(k)$ 表达式中的 $\hbar^2/(ma^2)$ 整体系数与后面的无量纲括号。
\end{itemize}
\end{warnbox}
